% =====================================================================
%  ANEXO C. Registro de ejecución de pruebas
% =====================================================================
\section{Registro de ejecución de pruebas}
\label{anx:ejecuciones}

\begin{notanormativa}[Registro nuevo respecto de la plantilla original]
Este registro corresponde a la actividad \textbf{TE3 «Registrar la ejecución de la prueba»} de
\normap{2}, y es independiente de la especificación del caso:

\begin{itemize}
  \item El \textbf{caso de prueba} (\ref{anx:casos}) describe \emph{qué debería ejecutarse}.
  \item La \textbf{ejecución} registra \emph{qué ocurrió una vez determinada}: en qué versión,
        en qué entorno, con qué datos, quién la realizó y con qué veredicto.
\end{itemize}

Un mismo caso puede ejecutarse muchas veces —en distintas versiones, tras una corrección, en
una campaña de regresión—. Cada una es una \textbf{ejecución distinta} con su propio
identificador. \textbf{Nunca sobrescriba el diseño del caso con el resultado de una ejecución}:
al hacerlo se pierde el historial y deja de poder demostrarse cuándo se corrigió un defecto.
\end{notanormativa}

\begin{instruccion}
Use la ficha detallada para las ejecuciones que requieran documentación amplia (típicamente las
no aprobadas) y la tabla resumen para el conjunto. Los campos de versión, configuración y
entorno son los que hacen reproducible el resultado: tómelos del Cuadro~\ref{tab:configuracion}.
\end{instruccion}

\subsection*{Ficha de ejecución}

\begin{xltabular}{\textwidth}{|E{4.4cm}|Y|}
\caption{Formato de registro de una ejecución de prueba.}\label{tab:ficha-ejecucion}\\ \hline
\endfirsthead
\multicolumn{2}{@{}l@{}}{\footnotesize\itshape\tablename~\thetable{} (continuación)}\\ \hline
\endhead
\multicolumn{2}{@{}r@{}}{\footnotesize\itshape continúa en la página siguiente}\\
\endfoot
\endlastfoot
  Identificador de ejecución      & \campo{EJ-014} \\ \hline
  Caso o procedimiento ejecutado  & \campo{CP-001 (v1) / PP-001 (v1)} \\ \hline
  Ciclo o campaña                 & \campo{Ciclo 2 — regresión} \\ \hline
  Versión del software probado    & \campo{v1.2.1 — identificación exacta de la compilación} \\ \hline
  Configuración utilizada         & \campo{Parámetros, banderas, navegador, sistema operativo} \\ \hline
  Entorno                         & \campo{ENT-xx y estado verificado (Anexo~F)} \\ \hline
  Datos utilizados                & \campo{DAT-xx, versión del conjunto} \\ \hline
  Fecha y hora                    & \campo{dd/mm/aaaa hh:mm} \\ \hline
  Ejecutor                        & \campo{Nombre} \\ \hline
  Entradas efectivamente usadas   & \campo{Valores concretos, si difieren de los especificados} \\ \hline
  Resultado esperado (referencia) & \campo{Referencia al campo del caso; no se transcribe} \\ \hline
  Resultado real observado        & \campo{Qué ocurrió realmente} \\ \hline
  Veredicto                       & \campo{Aprobado / no aprobado / no concluyente / bloqueado} \\ \hline
  Evidencia                       & \campo{EVID-014 — ruta, captura, registro o archivo} \\ \hline
  Incidente asociado              & \campo{INC-05, si el veredicto no fue aprobado} \\ \hline
  Observaciones                   & \campo{Anomalías del entorno, desviaciones del procedimiento} \\ \hline
\end{xltabular}
\notatabla{Veredicto \textit{no concluyente}: no pudo determinarse si el comportamiento fue correcto. \textit{Bloqueado}: la ejecución no pudo completarse por una causa externa al objeto de prueba.}

El Cuadro~\ref{tab:ficha-ejecucion} conserva, para cada intento, las condiciones exactas en que
se obtuvo el veredicto. Es el eslabón que une el caso diseñado con la evidencia y el incidente.

\subsection*{Registro consolidado de ejecuciones}

\begin{longtable}{|C{1.5cm}|C{1.6cm}|C{1.5cm}|C{1.7cm}|C{1.5cm}|C{1.9cm}|C{1.5cm}|C{1.5cm}|}
\caption{Registro consolidado de ejecuciones de prueba.}
\label{tab:registro-ejecuciones} \\
\hline
\filaenc \textbf{ID ejec.} & \textbf{Caso/proc.} & \textbf{Versión SW} & \textbf{Entorno} & \textbf{Fecha} & \textbf{Ejecutor} & \textbf{Veredicto} & \textbf{Incidente} \\ \hline
\endfirsthead
\hline
\filaenc \textbf{ID ejec.} & \textbf{Caso/proc.} & \textbf{Versión SW} & \textbf{Entorno} & \textbf{Fecha} & \textbf{Ejecutor} & \textbf{Veredicto} & \textbf{Incidente} \\ \hline
\endhead
\hline \multicolumn{8}{|r|}{\footnotesize\textit{continúa en la página siguiente}} \\ \hline
\endfoot
\hline
\endlastfoot
\campo{EJ-001} & \campo{CP-001} & \campo{v1.2.0} & \campo{ENT-01} & \campo{dd/mm} & \campo{Nombre} & \campo{No aprob.} & \campo{INC-05} \\ \hline
\campo{EJ-014} & \campo{CP-001} & \campo{v1.2.1} & \campo{ENT-01} & \campo{dd/mm} & \campo{Nombre} & \campo{Aprobado} & \campo{INC-05 (conf.)} \\ \hline
 & & & & & & & \\ \hline
 & & & & & & & \\ \hline
\end{longtable}

El Cuadro~\ref{tab:registro-ejecuciones} muestra el patrón esperado: el mismo caso CP-001
aparece dos veces, con identificadores de ejecución distintos y versiones distintas del
software. La segunda es la \emph{prueba de confirmación} del incidente INC-05
(apartado~\ref{sec:confirmacion-regresion}); la primera no se borra ni se modifica.
